% D04-003. AI-01 P0095--P0097; численные рейтинги не переносятся.
\section{Обучение и порядок сравнения архитектур}
\label{sec:ch4_training_protocol}

\workinprogress{Сохранён предварительный фрагмент протокола первой статьи. Раздел ещё не объединяет различающиеся опыты, обучение и разбиения данных остальных работ; готовым описанием методики он не является.}


\iffalse
(Примечание автора: Как-то совсем ни о чём, прост набросок, почти пустой. Ведь в разных статьях были разные исследования, и надо бы их аккуратно притянуть, а тут прост общее обучение, которое даже к методике ближе как будто. В общем совсем плохо глава и не о чём. Если это была стартовая попытка, то говори и прям в ПДФ итоговом, и/или делай ПДФ с комментариями, чтобы можно было видеть это, мои комментарии можно помечать как закрытые, если выполнены и я в них пометил «Ок», или что-то такое, без этого, даже если ты их решила, они висящие ещё. И всё равно, про случаи, когда глава является лишь наброском, прям в ПДФ должно быть сказано об этом, ибо вот к этой – очень сильные претензии) **[AA070: отложено; замечание открыто]**
(Примечание автора: Замечание закрыть, достаточно фразы о том, что процесс идёт и вычитка не требуется.)
\fi
\authorreviewnote{AA070}


Для сравнения архитектур используется задача распознавания
событий в реальных осциллограммах. В архитектурном исследовании
рассматривались нормальные режимы, рабочие коммутации,
аномальные режимы и аварийные события. Такое разделение
позволяет проверять различение повреждений и неаварийных
переходных процессов. Правило получения целевой метки из
временной разметки задаётся вместе с составом входного окна.

В исследовании архитектур обучение выполнялось в течение
60 эпох. Использовалась функция потерь BCEWithLogitsLoss,
которая получает выходные оценки до сигмоидального
преобразования. Веса учитывали неравномерную представленность
классов. Это связывает настройку модели с редкими событиями,
которые могут слабо влиять на общую долю правильных ответов.
При этом веса обучения и способ усреднения проверочных метрик
являются разными параметрами процедуры.

Из записей выбирались фрагменты для обучения. Применялись
масштабирование сигналов, изменение полярности и перестановка
фаз. Такие преобразования расширяют разнообразие входных
примеров. Для каждого преобразования требуется согласование
с целевой меткой. В частности, инвариантность признака наличия
повреждения не означает инвариантности требуемого направления
выхода РНМ.

Сравниваются состояние модели с наибольшим F1 на валидационной
выборке и состояние после последней эпохи. Первое показывает
результат выбора в ходе обучения, второе позволяет оценивать
изменение качества к его завершению. Проверочная выборка
используется для оценки выбранного состояния. Далее эти
результаты рассматриваются совместно с числом параметров,
временем расчёта и характером ошибок в отдельных режимах.
